% ==================== PARTIE F — QUESTIONS SUPPLÉMENTAIRES ====================
% Questions plus difficiles déplacées en fin de sujet (retours du référent) :
%   F1  QCM — précision flottante (ex-A2)
%   F2  QCM — évaluation court-circuit de || (ex-B1)

\element{var-qcm-flottant}{
    \begin{question}{var-qcm-flottant}
        La borne cumule trois tranches de facturation de 0,1 € chacune, puis vérifie le total.

        \begin{center}
            \begin{minipage}{.72\textwidth}
                \lstinputlisting[language=c]{sources/codes/a2-qcm-flottant.c}
            \end{minipage}
        \end{center}

        Qu'affiche ce fragment, et pour quelle raison ?
        \begin{reponses}
            \bonne{\texttt{ECART} : 0.1 n'a pas de représentation binaire exacte, le cumul des arrondis rend le total différent de 0.3}
            \mauvaise{\texttt{EXACT} : trois additions de 0.1 donnent exactement 0.3}
            \mauvaise{\texttt{ECART} : un \lstinline[language=C]|float| ne conserve que deux décimales}
            \mauvaise{Rien : comparer deux \lstinline[language=C]|float| avec \lstinline[language=C]|==| provoque une erreur de compilation}
        \end{reponses}
    \end{question}
}

\element{cond-qcm-courtcircuit}{
    \begin{question}{cond-qcm-courtcircuit}
        La fonction de lecture du capteur de tension est lente : chaque lecture est tracée
        à l'écran. Le module de sécurité est écrit ainsi :

        \begin{center}
            \begin{minipage}{.80\textwidth}
                \lstinputlisting[language=c]{sources/codes/b1-court-circuit.c}
            \end{minipage}
        \end{center}

        Qu'affiche exactement ce programme ?
        \begin{reponses}
            \bonne{\texttt{Coupure} seul : l'opérateur \lstinline[language=C]!||! n'évalue pas son second opérande lorsque le premier est vrai}
            \mauvaise{\texttt{[lecture capteur]} puis \texttt{Coupure} : les deux opérandes d'une condition sont toujours évalués}
            \mauvaise{\texttt{[lecture capteur]} seul : la condition est fausse car la tension dépasse la plage autorisée}
            \mauvaise{\texttt{Coupure} seul : un \lstinline[language=C]|printf| placé dans une fonction appelée depuis un \lstinline[language=C]|if| n'est jamais exécuté}
        \end{reponses}
    \end{question}
}
